\documentclass[a4paper,o12pt]{article}
\usepackage{librebaskerville}
\usepackage{xcolor}
\usepackage{ragged2e}

\definecolor{darkblue}{HTML}{16247C}
\definecolor{black}{HTML}{000000}
%----------------------------------------------------------------------------------------
%	FONT
%----------------------------------------------------------------------------------------

%----------------------------------------------------------------------------------------
%	PACKAGES
%----------------------------------------------------------------------------------------
\usepackage{url}
\usepackage{parskip} 	

%other packages for formatting
\RequirePackage{color}
\RequirePackage{graphicx}
\usepackage[usenames,dvipsnames]{xcolor}
\usepackage[scale=0.9]{geometry}

%tabularx environment
\usepackage{tabularx}

%for lists within experience section
\usepackage{enumitem}

% centered version of 'X' col. type
\newcolumntype{C}{>{\centering\arraybackslash}X} 

%to prevent spillover of tabular into next pages
\usepackage{supertabular}
\usepackage{tabularx}
\newlength{\fullcollw}
\setlength{\fullcollw}{0.47\textwidth}

%custom \section
\usepackage{titlesec}				
\usepackage{multicol}
\usepackage{multirow}

%CV Sections inspired by: 
%http://stefano.italians.nl/archives/26
\titleformat{\section}{\Large\scshape\raggedright}{}{0em}{}[\titlerule]
\titlespacing{\section}{0pt}{10pt}{10pt}

%for publications
\usepackage[style=authoryear,sorting=ynt, maxbibnames=2]{biblatex}

%Setup hyperref package, and colours for links
\usepackage[unicode, draft=false]{hyperref}
\definecolor{linkcolour}{rgb}{0,0.2,0.6}
\hypersetup{colorlinks,breaklinks,urlcolor=linkcolour,linkcolor=linkcolour}
\addbibresource{citations.bib}
\setlength\bibitemsep{1em}

%for social icons
\usepackage{fontawesome5}

%debug page outer frames
%\usepackage{showframe}


% job listing environments
\newenvironment{jobshort}[2]
    {
    \begin{tabularx}{\linewidth}{@{}l X r@{}}
    \textbf{#1} & \hfill &  #2 \\[3.75pt]
    \end{tabularx}
    }
    {
    }

\newenvironment{joblong}[2]
    {
    \begin{tabularx}{\linewidth}{@{}l X r@{}}
    \textbf{#1} & \hfill &  #2 \\[3.75pt]
    \end{tabularx}
    \begin{minipage}[t]{\linewidth}
    \begin{itemize}[nosep,after=\strut, leftmargin=1em, itemsep=3pt,label=--]
    }
    {
    \end{itemize}
    \end{minipage}    
    }



%----------------------------------------------------------------------------------------
%	BEGIN DOCUMENT
%----------------------------------------------------------------------------------------
\begin{document}

% non-numbered pages
\pagestyle{empty} 

%----------------------------------------------------------------------------------------
%	TITLE
%----------------------------------------------------------------------------------------

% \begin{tabularx}{\linewidth}{ @{}X X@{} }
% \huge{Your Name}\vspace{2pt} & \hfill \emoji{incoming-envelope} email@email.com \\
% \raisebox{-0.05\height}\faGithub\ username \ | \
% \raisebox{-0.00\height}\faLinkedin\ username \ | \ \raisebox{-0.05\height}\faGlobe \ mysite.com  & \hfill \emoji{calling} number
% \end{tabularx}

\begin{tabularx}{\linewidth}{@{} C @{}}
\begin{flushleft}
	\fontsize{30}{60}\selectfont{{\textbf{\color{darkblue}Ícaro Onofre Silva.}}} \\[10pt]
\end{flushleft}
\color{black}
\href{https://github.com/icaro-onofre}{\raisebox{-0.05\height}\faGithub\ icaro-onofre} \ $|$ \ 
\href{https://linkedin.com/in/icaro-onofre}{\raisebox{-0.05\height}\faLinkedin\ icaro-onofre} \ $|$ \ 
\href{mailto:icaro.onofre.s@gmail.com}{\raisebox{-0.05\height}\faEnvelope \ icaro.onofre.s@gmail.com} \ $|$ \ 
\href{tel:+11937045276}{\raisebox{-0.05\height}\faMobile \ +55 11937045276} \\
\end{tabularx}
%----------------------------------------------------------------------------------------
%	EDUCATION
%----------------------------------------------------------------------------------------
\section{{\textcolor{darkblue}{\textbf{Educação}}}}
\begin{tabularx}{\linewidth}{@{}l X@{}}	

	2017 - 2019 Ensino médio integrado a técnico em informática para internet \textbf{Etec CLP.} \\

2020 - 2024 Análise e desenvolvimento de sistemas - \textbf{Fatec Dep Ary Fossen de Jundiaí} \\

\end{tabularx}

\section{{\textcolor{darkblue}{\textbf{Certificados}}}}
\begin{tabularx}{\linewidth}{@{}l X@{}}	

Dezembro 2020 - Python para análise de dados  - \textbf{Impacta} \\

Junho 2021 - ReactJS - \textbf{Alura} \\

\end{tabularx}


%----------------------------------------------------------------------------------------
% EXPERIENCE SECTIONS
%----------------------------------------------------------------------------------------
%Experience
\section{\textcolor{darkblue}{\textbf{Experiência}}}
\begin{joblong}{Desenvolvedor de software Júnior (Khanum)}{Junho 2022 - Janeiro 2023}
\item Eu desenvolvi o frontend de um aplicativo para rede social que ajudava conectar médicos estetas a pacientes com React Native e ReactJS
    \begin{enumerate}
    \item Criação de telas para cadastro, login, feed de postagens e mensagens.
    \end{enumerate}
\end{joblong}

\begin{joblong}{Estágio (LAZco)}{Agosto 2021 - Dezembro 2021}
\item Eu realizei manutenções num ERP legado desenvolvido em Delphi 7, adicionando novas funcionalidades a telas de gestão.
	\begin{enumerate}
		\item Interface para gestão de dependentes num clube de lazer.
		\item Cadastro de CIDs de acidentes de trabalho.
		\item Treinei um novo desenvolvedor que estava ingressando na empresa.
	\end{enumerate}
	
\end{joblong}

\begin{joblong}{Unifrota}{Julho 2023 - Presente}
\item Eu desenvolvi novas novas funcionalidades para um ERP legado, um aplicativo e um dashboard web para gestão de frotas de veiculos, e também contribui para 
um sistema de checking de OOH
	\begin{enumerate}
		\item Ponto eletronico para gestão de funcionários.
		\item Controle das horas de trabalho da jornada de veículos.
		\item Um sistema que rastreeava a vida útil de pneus e controlava a frequência que eles deviam ser recapados e ter seus sulcos medidos.
		\item Ajudei a otimizar os custos de computação em nuvem  ao remover um sistema legado ja obsoleto.
		\item Relatórios de prova de postagem de campanhas em painéis digitais.
	\end{enumerate}
	
\end{joblong}


\section{\textcolor{darkblue}{\textbf{ Projetos }}}


\begin{tabularx}{\linewidth}{ @{}l r@{} }
\textbf{Okonomicus} & \hfill \href{https://some-link.com}{www.okonomicus.com} \\[3.75pt]
\multicolumn{2}{@{}X@{}}{ Okonomicus é um sistema de análise fundamentalista de ações que estou desenvolvendo com meu pai, meu principal
	interesse é compilar tudo que eu valorizo quando analiso ações e gerencio aspectos financeiros de um projeto de ambição pessoal, e aprimorar meus conhecimentos 
    do mercado financeiro.}  \\
    \\

\textbf{Studio de Idiomas} & \hfill \href{http://studiodeidiomas.com.br/}{studiodeidiomas.com.br} \\[3.75pt]
\multicolumn{2}{@{}X@{}}{ Realizei a renovação de uma landing page para uma escola de idiomas, criando um design mais moderno e mais intuitivo. }  \\

\end{tabularx}

% \begin{tabularx}{\linewidth}{ @{}l r@{} }
% \end{tabularx}

\vfill
\center{\footnotesize Last updated: \today}

\end{document}

